% Replaces "Synthetic ground-truth accuracy" in pocketdrs_paper.tex.
% Source: server/scripts/synth_validate.py --wide, 100 deliveries, ground truth from the
% generated trajectory. Raw: dump/validation/wide_grid_summary.csv
%
% NIRAJ: read these before using them. The headline drops from 8/8 to 64/100 and the paper is
% far stronger for it. Decide whether you agree with the reading.

\subsection{Verdict accuracy}
We evaluate on 100 synthetic deliveries spanning five release speeds (90 to 145\,km/h), five
lines (middle, and off and leg at 20 and 40\,cm) and four lengths (2 to 8\,m). Ground truth is
taken from the generating trajectory under the same geometry the decision layer uses. The set is
82 not out and 18 out, so a policy that answers one class alone cannot score well.

PocketDRS agreed with the ground-truth verdict in \textbf{64 of 100} deliveries. It returned no
verdict on 10, where calibration or the arc fit failed outright; agreement over the 90 that
produced one is 71\,\%.

\begin{table}[!t]
\caption{Verdict outcomes over the 100-delivery grid.}
\label{tab:wide}
\centering
\begin{tabular}{@{}llr@{}}
\toprule
\textbf{Ground truth} & \textbf{Predicted} & \textbf{n}\\
\midrule
not out & not out       & 48\\
out     & out           & 16\\
not out & umpire's call & 19\\
not out & out           & \textbf{7}\\
not out & no verdict    & 8\\
out     & no verdict    & 2\\
\bottomrule
\end{tabular}
\end{table}

Two features of Table~\ref{tab:wide} matter more than the headline. The 19 not-out deliveries
returned as umpire's call are the intended behaviour of the widened bands: the system declines to
assert rather than guessing, and a coaching tool that says "too close to call" is useful. The
seven not-out deliveries returned as \textsc{out} are not. A false out is the error a review
system must not make, and it is the one number a deployment decision should turn on.

Accuracy is strongly structured, and the structure is physical rather than incidental:

\begin{table}[!t]
\caption{Verdict agreement by delivery parameter.}
\label{tab:wide_break}
\centering
\begin{tabular}{@{}llr@{}}
\toprule
\textbf{Parameter} & \textbf{Value} & \textbf{Agreement}\\
\midrule
Length & 6--8\,m (good, short) & 76\,\%\\
       & 4\,m (full)           & 56\,\%\\
       & 2\,m (yorker)         & 48\,\%\\
\midrule
Line   & middle stump          & 80\,\%\\
       & $\pm$40\,cm           & 70--80\,\%\\
       & $\pm$20\,cm           & 45\,\%\\
\midrule
Speed  & 90--105\,km/h         & 70\,\%\\
       & 130--145\,km/h        & 60\,\%\\
\bottomrule
\end{tabular}
\end{table}

A fuller length gives a longer post-bounce arc for the solver to fit, so accuracy rises with
length. The dip at $\pm$20\,cm is the marginal band: those deliveries pass close to the outer
stump, where the depth error quantified in Section~\ref{sec:setup} is enough to move the predicted
crossing across the decision boundary. Deliveries well outside the line are easy again, which is
why $\pm$40\,cm recovers. Faster deliveries give fewer frames between bounce and stump plane.

This supersedes an earlier evaluation on eight deliveries at a single speed, one length band and
lines within $\pm$6\,cm of middle stump, which reported full agreement. Every delivery in that set
was out and every one sat in the region Table~\ref{tab:wide_break} identifies as the system's
strongest, so the figure measured neither discrimination nor range. We report the wider grid
instead, including where it fails.
